%% ============================================================================
%%  chapters/chapter2.tex
%%  فصل ۲: پیشینه پژوهش و مبانی نظری
%%  قالب رسمی پایان‌نامه و رساله دانشگاه خوارزمی
%% ============================================================================

\chapter{پیشینه پژوهش و مبانی نظری}
\label{chap:literature}

\section{مقدمه}
\label{sec:ch2_intro}

درک عمیق مسائل مطرح در حروف‌چینی دیجیتال، بازنمایی ساختار اسناد، و خودکارسازی صفحه‌آرایی نیازمند بررسی پایه‌های نظری در نظریه ارتباطات، سامانه‌های محاسباتی، تایپوگرافی رایانه‌ای، و رویکردهای نوین یادگیری ماشین است. در این فصل، ابتدا مبانی نظری حاکم بر پردازش اطلاعات، تکوین موتورهای حروف‌چینی دوسویه، استانداردهای بین‌المللی نویسه‌ها، و الگوهای یادگیری عمیق مورد واکاوی قرار می‌گیرد. سپس، پژوهش‌های پیشین انجام‌شده در سطح بین‌المللی و مطالعات پیشگام پژوهشگران داخلی در حوزه صفحه‌آرایی زبان فارسی و پردازش زبان طبیعی بررسی و در قالب یک جدول تطبیقی جامع مقایسه می‌گردند.

\section{مبانی نظری تحقیق}
\label{sec:ch2_theory}

\subsection{نظریه اطلاعات و بنیان‌های ارتباطی}
پایه‌های بنیادین انتقال و کدگذاری داده‌ها در سده بیستم توسط کلود شانون\LTRfootnote{\lr{Claude Elwood Shannon}} پی‌ریزی گردید. شانون در اثر تاریخی خود با عنوان نظریه ریاضی ارتباطات\index{نظریه اطلاعات} \cite{shannon1948mathematical}، مفهوم آنتروپی اطلاعات\index{آنتروپی اطلاعات} را معرفی کرد و نشان داد که محتوای پیام را می‌توان به صورت کمّی اندازه‌گیری و بدون اتلاف در کانال‌های دارای نوفه منتقل نمود. از سوی دیگر، آلن تورینگ\LTRfootnote{\lr{Alan Mathison Turing}} با طرح پرسش اساسی در خصوص قابلیت تفکر ماشین‌ها \cite{turing1950computing}، افق‌های نوینی را در پردازش نمادین و منطق محاسباتی گشود که بعدها به ظهور هوش مصنوعی\index{هوش مصنوعی} مدرن انجامید. این دو دیدگاه زیربنای نظری مدل‌سازی نمادها و متون دیجیتال را شکل می‌دهند.

\subsection{تک، لاتک و تکوین تایپوگرافی دیجیتال}
تا پیش از اواخر دهه ۱۹۷۰ میلادی، چاپ فرمول‌های ریاضی و متون علمی نیازمند فرآیندهای دستی پرهزینه و پرخطا در چاپخانه‌ها بود. دانلد کنوت\LTRfootnote{\lr{Donald Ervin Knuth}} با هدف ایجاد سامانه حروف‌چینی دقیق و مستقل از دستگاه، سامانه تک را ابداع نمود و اصول آن را در کتاب مشهور تک‌بوک \cite{knuth1984texbook} به تفصیل مدون کرد. نوآوری کنوت در ابداع الگوریتم شکست پاراگراف به سطرها بر اساس مفهوم کمینه‌سازی جریمه\LTRfootnote{\lr{Penalty and Badness Optimization}}، هنوز هم استاندارد طلایی صفحه‌آرایی متون در جهان به شمار می‌آید. پس از آن، لزلی لمپورت با توسعه بسته ماکروی لاتک، تمرکز نویسنده را از جزئیات مکانیکی صفحه به ساختار منطقی سند معطوف کرد.

\subsection{الگوریتم دوسویه یونیکد و موتور زی‌پرشین}
خط فارسی برخلاف خطوط لاتین از راست به چپ نوشته می‌شود، در حالی که ارقام، علائم جبری و عبارات لاتین تعبیه‌شده در آن جهت چپ‌به‌راست دارند. کنسرسیوم یونیکد \cite{unicode2024standard} با تدوین پیوست استاندارد شماره ۹ (\lr{UAX \#9})، الگوریتم حل جهت دوسویه\index{متن دوسویه}\LTRfootnote{\lr{Unicode Bidirectional Algorithm (Bidi)}} را معرفی نمود. وفا خلیقی \cite{khalighi2020xepersian} با تکیه بر قابلیت‌های موتور مدرن زی‌تک، بسته زی‌پرشین را پیاده‌سازی کرد. این بسته با بازتعریف پارادایم‌های پایه لاتک، پشتیبانی بی‌نقصی از ویژگی‌های ذاتی خط و زبان فارسی نظیر اتصال حروف، تنوع قلم‌های نستعلیق و نسخ، کشیدگی کلمات و تنظیم خودکار پانویس‌ها را به ارمغان آورد.

\subsection{یادگیری ماشین و شبکه‌های عصبی عمیق}
در دهه گذشته، پارادایم پردازش اسناد و تحلیل متون با پیشرفت‌های شگرف یادگیری ماشین\index{یادگیری ماشین} دگرگون شده است. مبانی یادگیری آماری و شناسایی الگو توسط بیشف \cite{bishop2006pattern} و اصول یادگیری عمیق\index{یادگیری عمیق} توسط گودفلو و همکاران \cite{goodfellow2016deep} تدوین شده است. در جامعه علمی کشور نیز کتاب مبانی یادگیری ماشین کریمی \cite{karimi1398} و ترجمه فارسی مرجع جامع هوش مصنوعی راسل و نورویگ توسط سلیمی \cite{salimi1400translation} به تبیین دقیق این مفاهیم برای دانشجویان و پژوهشگران پرداخته‌اند.

علاوه بر این، معرفی معماری ترنسفورمر\index{ترنسفورمر} و سازوکار توجه خودکار\index{سازوکار توجه} توسط واسوانی و همکاران \cite{vaswani2017attention}، انقلابی در پردازش توالی‌ها پدید آورد که منجر به پیدایش مدل‌های زبانی عظیم گردید. در حوزه پردازش تصویر و بازشناسی نویسه‌ها، شبکه‌های زاینده تخاصمی\index{شبکه زاینده تخاصمی} گودفلو و همکاران \cite{goodfellow2014generative} و شبکه‌های عمیق بازمانده\index{شبکه عصبی پیچشی} هه و همکاران \cite{he2016deep} بن‌بست‌های یادگیری در شبکه‌های بسیار عمیق را مرتفع ساختند.

\section{مرور پژوهش‌های پیشین}
\label{sec:ch2_related_works}

\subsection{پژوهش‌های بین‌المللی در خودکارسازی صفحه‌آرایی اسناد}
تحقیقات در زمینه درک ساختار اسناد\LTRfootnote{\lr{Document Layout Analysis}} و تولید خودکار متون دانشگاهی طی سال‌های اخیر شتاب فزاینده‌ای یافته است. محققان دانشگاه‌های استنفورد و ام‌آی‌تی با توسعه بسته‌های متعددی برای مجلات معتبر \lr{IEEE} و \lr{ACM}، استانداردسازی خروجی مقالات را محقق ساختند. تمرکز اصلی این پژوهش‌ها بر تفکیک محتوا از قالب، کاهش بار محاسباتی کامپایلر، و تولید اسناد دسترسی‌پذیر\LTRfootnote{\lr{Accessible PDF}} بوده است. با این حال، غالب این سامانه‌ها برای زبان‌های تک‌جهته طراحی شده و در تعامل با خطوط عربی و فارسی با محدودیت‌های اساسی روبه‌رو هستند.

\subsection{پژوهش‌های داخلی در نگارش متون دانشگاهی و پردازش خط فارسی}
در داخل کشور، تلاش‌های متعددی برای استانداردسازی حروف‌چینی و پردازش متون علمی صورت پذیرفته است:
\begin{itemize}
  \item \textbf{مطالعه سامانه‌های اطلاعاتی و رفتار پژوهشگران:} رجبعلی‌بگلو، فتاحی و پریرخ \cite{fatahi1395} به بررسی تأثیر ساختار نظام‌های اطلاعاتی و قالب‌های نگارشی بر رفتار اطلاع‌یابی و بهره‌وری پژوهشگران پرداختند و نشان دادند که استانداردسازی فرمت اسناد، سرعت بازیابی اطلاعات و ارجاع‌دهی را به میزان چشمگیری افزایش می‌دهد.
  \item \textbf{مدل‌سازی زبان فارسی با شبکه‌های عمیق:} دهقان \cite{dehghan1399phd} در رساله دکترای خود در دانشگاه خوارزمی به بررسی بهینه‌سازی شبکه‌های عصبی بازگشتی و ترنسفورمرها در پردازش زبان طبیعی فارسی پرداخت و چالش‌های واژه‌نگاشت، فاصله‌گذاری مجازی و ساختار دستوری متون علمی فارسی را مورد ارزیابی قرار داد.
  \item \textbf{بهینه‌سازی توزیع‌شده در یادگیری عمیق:} شمس، محمدی و رضایی \cite{shams1401} الگوریتم‌های بهینه‌سازی توزیع‌شده را برای آموزش مدل‌های عمیق بررسی نمودند که نتایج آن در تسریع زمان آموزش سامانه‌های هوشمند پردازش سند بسیار راه‌گشا است.
  \item \textbf{صفحه‌آرایی خودکار در زبان‌های راست‌به‌چپ:} احمدی \cite{ahmadi1401thesis} در پایان‌نامه کارشناسی ارشد خود در دانشکده مهندسی برق و کامپیوتر دانشگاه خوارزمی، مسئله طراحی سامانه‌های خودکار صفحه‌آرایی اسناد دانشگاهی را در بستر زبان‌های راست‌به‌چپ تحلیل و ساختار کلاس‌های تک را بررسی نمود.
  \item \textbf{تحلیل مقایسه‌ای سامانه‌های حروف‌چینی متن‌باز:} رضایی و حسینی \cite{rezaei1402} در کنفرانس ملی کامپیوتر ایران، عملکرد سامانه‌های مختلف حروف‌چینی را در تولید مستندات علمی دوزبانه مقایسه کردند و نشان دادند که زی‌پرشین در مقایسه با سایر ابزارها دارای بالاترین امتیاز در حفظ زیبایی بصری و انطباق با ضوابط دانشگاهی است.
\end{itemize}

\section{تحلیل مقایسه‌ای و جدول تطبیقی کارهای مرتبط}
\label{sec:ch2_comparison}

به منظور ارزیابی جامع رویکردهای موجود و تبیین وجوه تمایز قالب ارائه‌شده در این پژوهش با سایر ابزارها، جدول \ref{tab:comparison_literature} ویژگی‌های کلیدی سامانه‌های متداول را مقایسه می‌نماید.

\begin{table}[htbp]
\centering
\caption{مقایسه تطبیقی سامانه‌های حروف‌چینی و قالب‌های اسناد دانشگاهی}
\label{tab:comparison_literature}
\vspace{4pt}
\small
\setlength{\tabcolsep}{4pt}
\resizebox{\linewidth}{!}{%
\begin{tabular}{|p{2.3cm}|c|c|c|c|p{4.2cm}|}
\hline
\textbf{سامانه / رویکرد} & \textbf{محیط اجرا} & \textbf{پشتیبانی دوزبانه} & \textbf{پایداری ارجاعات} & \textbf{انطباق با خوارزمی} & \textbf{مهم‌ترین مزایا و محدودیت‌ها} \\
\hline
مایکروسافت ورد & تجاری & متوسط & ضعیف & اسمی (فایل الگو) & سهولت کاربری اولیه، اما به‌هم‌ریختگی در اسناد حجیم و فرمول‌ها \\
\hline
لاتک استاندارد & متن‌باز & نامناسب (نیازمند عرب‌تک) & عالی & ندارد & توانایی بالا در متون انگلیسی، ناتوانی در قلم‌های بومی فارسی \\
\hline
قالب‌های متفرقه وب & متن‌باز & ناپایدار & متوسط & ندارد & عدم هماهنگی با حاشیه‌ها و سرصفحه‌های شیوه‌نامه مصوب ۱۴۰۲ \\
\hline
قالب دانشگاه تهران & متن‌باز & خوب & خوب & مغایر در حاشیه‌ها & تنظیم حاشیه‌ها بر اساس ۳/۵ و ۲/۵ سانتی‌متر (ناهمخوان با خوارزمی) \\
\hline
\textbf{قالب پیشنهادی خوارزمی} & \textbf{متن‌باز} & \textbf{کامل و پایدار} & \textbf{عالی (دوزبانه)} & \textbf{۱۰۰ درصد منطبق} & \textbf{طراحی ماژولار، قلم‌های رسمی بی زر و بی تیتر، حذف کامل هشدارهای نویسه} \\
\hline
\end{tabular}%
}
\end{table}

\section{جمع‌بندی فصل}
\label{sec:ch2_summary}

بررسی مبانی نظری و پیشینه پژوهش‌های پیشین نشان می‌دهد که حروف‌چینی اسناد علمی فارسی در سطوح عالی دانشگاهی، نیازمند اتخاذ رویکردی چندبعدی شامل تسلط بر استانداردهای یونیکد، قواعد دوسویه، معماری ماژولار لاتک و تطبیق ریاضی با فواصل شیوه‌نامه دانشگاه خوارزمی است. اگرچه تلاش‌های ارزنده‌ای توسط پژوهشگران گذشته صورت گرفته، اما خلأ یک قالب یکپارچه، کاملاً آزموده‌شده، فاقد هرگونه خطای نویسه گم‌شده، و همراه با مستندات کامل برای دانشجویان دانشگاه خوارزمی همچنان احساس می‌شد. مدل پیشنهادی که در فصل بعدی تشریح می‌گردد، پاسخگوی دقیق این نیازمندی‌ها خواهد بود.
